\unit{Quadratures and phase-space pictures}
\unitref{W1L3 and W3L1; homodyne lectures}
\subsection{What are the two quadratures of a wave?}
A monochromatic classical electric field can be written
\[
E(t)=A\cos\omega t+B\sin\omega t.
\]
The coefficients \(A,B\) specify its amplitude and
phase. They are not frequency and amplitude: the
\emph{frequency} \(\omega\) is already fixed by
the mode. Quantum mechanically, define
\[
\boxed{x=a+a^\dagger,\qquad
p=i(a^\dagger-a),\qquad a=(x+ip)/2.}
\]
The commutator is
\[
[x,p]=[a+a^\dagger,i(a^\dagger-a)]
=2i.
\]
Heisenberg's inequality gives
\[
\boxed{\Delta x\,\Delta p\ge1.}
\]
In this course convention the vacuum has
\((\Delta x)^2=(\Delta p)^2=1\).
A textbook using \(x=(a+a^\dagger)/\sqrt2\)
has vacuum variance \(1/2\); no physics
changes, only the scaling.

\subsection{The meaning of the phase-space plane}
For a coherent state \(\ket\alpha\),
\[
\langle x\rangle=2\operatorname{Re}\alpha,\quad
\langle p\rangle=2\operatorname{Im}\alpha,\quad
(\Delta x)^2=(\Delta p)^2=1.
\]
Thus a coherent state looks like a
vacuum uncertainty circle displaced from the
origin. Its centre identifies a classical-like
field amplitude; its width encodes unavoidable
quantum fluctuations.

\begin{center}
\begin{tikzpicture}[>=Latex,scale=.8]
\draw[->](-2.2,0)--(6.2,0)node[right]{\(x\)};
\draw[->](0,-2.3)--(0,3.4)node[above]{\(p\)};
\draw[gray,dashed,thick](0,0)circle(1);
\draw[accent,fill=accent!10,thick](3.4,1.3)circle(1);
\fill[accent](3.4,1.3)circle(.07);
\draw[dotted](0,0)--(3.4,1.3);
\node[below left]at(-.4,-1){vacuum};
\node[above right]at(4.1,1.9){coherent};
\end{tikzpicture}
\end{center}

\subsection{Rotated quadratures}
For a local-oscillator phase \(\theta\),
\[
\boxed{X_\theta=ae^{-i\theta}
+a^\dagger e^{i\theta}
=x\cos\theta+p\sin\theta.}
\]
Homodyne detection measures the probability
distribution of \(X_\theta\). At \(\theta=0\)
it measures \(x\); at \(\theta=\pi/2\)
it measures \(p\). Measuring distributions
over many angles enables optical homodyne
tomography. The phase-space distribution
is reconstructed from different projections,
rather than measured pointwise with
arbitrarily sharp \(x,p\) simultaneously.

\begin{worked}{Rotate the local oscillator}
Take \(\alpha=1+i\). Then
\(\langle x\rangle=\langle p\rangle=2\).
At \(\theta=\pi/4\),
\[
\langle X_{\pi/4}\rangle
=2\cos(\pi/4)+2\sin(\pi/4)=2\sqrt2.
\]
Its variance remains 1, since coherent
vacuum noise is equally distributed
in all directions.
\end{worked}

\subsection{Three complementary quasiprobabilities}
The Wigner function \(W(x,p)\) represents
a quantum state in phase space. Its
marginals give measurable quadrature
probabilities, but some quantum states
have negative Wigner regions.
The Husimi function is
\[
Q(\alpha)=\frac1\pi\bra\alpha\rho\ket\alpha,
\]
which is always nonnegative and is smoother
than the Wigner function. The
Glauber--Sudarshan representation is
\[
\rho=\int P(\alpha)\ket\alpha\bra\alpha\,d^2\alpha.
\]
Here \(P\) can be singular or nonpositive:
it is \emph{not always} a classical
probability distribution.
All three describe the same \(\rho\), but
highlight different physical properties.

\subsection{Gaussian states and covariance}
For \(R=(x,p)^T\), define
\[
d_i=\langle R_i\rangle,\qquad
V_{ij}=\tfrac12\langle\Delta R_i\Delta R_j+
\Delta R_j\Delta R_i\rangle.
\]
A Gaussian state is completely specified
by \(d,V\). Coherent and squeezed Gaussian
states have Gaussian Wigner functions.
A lossless linear transformation \(R'=SR\)
acts as \(d'=Sd,\ V'=SVS^T\).
An open-system Gaussian channel generally adds
noise: \(V'=SVS^T+N\).

\begin{checkpoint}
Phase space is useful for \emph{states of modes}.
Quadrature \(x\) is not position of a photon
and \(p\) is not its momentum. They are
conjugate field-amplitude variables.
\end{checkpoint}
